\section{STRATIFIED RANDOM-ROLE CALIBRATION}
\label{app:stratified}
\paragraph{A remedy based on randomized roles.}
Rather than use a low-exposure reference for every test node, retain both exposure strata and calibrate each against its own randomly acquired normals. Fix the strata using the graph, a separate training panel, and observed anomalous neighbors. Their construction must be invariant to which remaining normal nodes become calibration or test nodes. A threshold estimated only from the realized calibration pool generally does not meet this requirement.

\begin{lemma}[Stratified random-role calibration]
\label{lem:stratified}
Condition on the graph, labels, deployment universe, training panel, role-invariant fitted scores, and all anomaly-side test and revelation assignments. Let $D_1,\ldots,D_H$ be strata fixed by this information. Resolve all score ties lexicographically with independent continuous marks drawn before normal-role assignment. Uniformly choose a prescribed number of deployment normals for testing, and independently reveal each remaining normal with probability $\rho$. In stratum $h$, let $n_h$ be the revealed-normal reference size, $m_{0h}$ the normal-test count, and $m_h$ the total test count. Apply conformal BH separately at fixed levels $\alpha_h$ with $\sum_h\alpha_h\le\alpha$, abstaining if the reference or test stratum is empty. Conditional also on all stratum role counts,
\begin{equation}
\E[\FDP_{\rm union}]
\le\sum_{h:m_h>0}\alpha_h\frac{m_{0h}}{m_h}\le\alpha.
\label{eq:stratified}
\end{equation}
\end{lemma}
The conditional expectation in Equation~\eqref{eq:stratified} averages over normal identities assigned to roles; it does not fix the realized normal test set. Uniform role assignment supplies the exchangeability needed within each stratum~\citep{marandon2024adaptive}. Splitting $\alpha$ avoids a cross-stratum dependence assumption because $\FDP_{\rm union}\le\sum_h\FDP_h$ pointwise. Small strata reduce rank resolution and power, while empty references safely abstain. The proof and the comparisons with pooled full-reference calibration follow.

\paragraph{Conditioning and construction.}
Let $\mathcal F$ include the finite graph, its entire fixed label vector, the training panel, the deployment universe, fitted scores unchanged by normal-role reassignment, and every anomalous node's test/revelation assignment. It also includes independent continuous score tie marks and any auxiliary marks used to break exposure ties, drawn before normal-role assignment. The partition $D_1,\ldots,D_H$ is $\mathcal F$-measurable. This conditioning is a design-based analysis; the procedure does not observe hidden test labels.

For example, the numerator of observed exposure counts anomalous neighbors known from the training and revealed non-test panels. Conditional on their assignments it is fixed; division by the graph's fixed degree preserves this property. Changing which normal neighbors are revealed contributes zero to the numerator. By contrast, dividing by the number of revealed neighbors, or choosing a threshold from just the realized normal reference, can make the partition depend on normal roles and requires another argument.

\paragraph{Proof of Lemma~\ref{lem:stratified}.}
Write $N$ for the number of deployment normals and $t$ for the prescribed normal-test count. For $0<\rho<1$, a particular normal-role partition into test nodes $T_0$, revealed reference nodes $C$, and unused nodes $U$ has probability
\[
\binom Nt^{-1}\rho^{|C|}(1-\rho)^{|U|}.
\]
This probability is constant over assignments having the same role counts in each fixed stratum. Thus, conditional on $\mathcal F$ and those counts, normal identities are exchangeable among the roles within each stratum. At $\rho=0$ or $1$, the same statement holds on the feasible assignments directly. Anomaly-side and normal-side sampling must have this conditional law; conditioning on an arbitrary coupled allocation would not establish it.

Fix a stratum and additionally condition on the unordered set of normals occupying calibration and null-test roles. Assign them to calibration and null-test positions, introducing independent random orderings within each group. Every permutation into those positions is equally likely. Hence their score vector is jointly exchangeable conditional on the anomalous test scores. Auxiliary ordering does not alter any p-value attached to a node or the FDP. Lexicographic tie marks make every calibration and test score distinct; equivalently use ranks in this fixed total ordering. The score-level exchangeability result of \citet[Theorem 3.3]{marandon2024adaptive} supplies null super-uniformity and PRDS within this stratum. BH therefore gives
\[
\E[\FDP_h\mid\mathcal F,\mathcal K]
\le\alpha_h m_{0h}/m_h,
\]
after averaging over the additional unordered-set conditioning, where $\mathcal K$ denotes all stratum role counts. Empty references abstain, and empty test strata have FDP zero.

Let $R_h$ and $V_h$ be the discovery and false-discovery counts. For every realization,
\[
\frac{\sum_h V_h}{(\sum_h R_h)\vee1}
\le\sum_h\frac{V_h}{R_h\vee1}.
\]
Taking conditional expectations proves Equation~\eqref{eq:stratified}. No independence or PRDS condition across strata is invoked. The strata and level allocation must be fixed under the stated conditioning; selecting them after inspecting FDP is not covered.

\paragraph{Prespecified empirical comparison.}
We use two separate binary partitions: zero versus positive observed exposure, and the lower 80\% versus upper 20\% of exposure ranks over the entire deployment universe. Independent continuous marks break exposure ties in the latter. Its cutoff is fixed under normal-role permutations, unlike a cutoff computed from only calibration normals. Neither partition is selected based on its results. Each stratum receives level $0.05$, including when another stratum is empty.

For each graph, both frozen supervised scorers, all ten training seeds, five test splits, and revealed fractions $0.25,0.50,1$, we compare pooled full-reference BH at $0.10$, low-stratum calibration applied to the entire test set at $0.10$, and separate-stratum BH. Full-reference and stratified calibration use the same total acquired labels and reference nodes. The low-stratum diagnostic uses fewer references; the separate dose experiment supplies size-matched composition comparisons. We retain every result, including power losses and abstentions.

\paragraph{Implementation checks.}
Enumeration over 44,256 joint role assignments checks both the stratum-level BH bound and the union bound; 16,110 assignments have nonempty discoveries. Two hundred tied-score examples compare rank-based BH with a direct reference implementation. Empty references abstain. For each evaluated split and revelation fraction, altering unknown labels or normal revelation roles leaves exposure unchanged; both partition constructions are checked for the same invariance. The protocol, per-split outcomes, seed summaries and input hashes are included in the supplement.

\begin{table*}[t]\centering\small
\caption{Amazon: complete stratified comparison, mean FDP / power over ten training-seed averages. Both partitions and all revelation fractions are shown; no partition is selected using these outcomes.}
\begin{tabular}{lllrrr}\toprule Score & $\rho$ & Partition & Low reference & Full random & Stratified\\\midrule
Attributes & 0.25 & Zero/positive & $0.093/0.679$ & $0.091/0.779$ & $0.039/0.567$ \\
Attributes & 0.25 & 80/20 ranks & $0.110/0.793$ & $0.091/0.779$ & $0.036/0.530$ \\
Attributes & 0.50 & Zero/positive & $0.140/0.770$ & $0.089/0.779$ & $0.042/0.651$ \\
Attributes & 0.50 & 80/20 ranks & $0.114/0.797$ & $0.089/0.779$ & $0.039/0.581$ \\
Attributes & 1.00 & Zero/positive & $0.198/0.807$ & $0.092/0.783$ & $0.044/0.694$ \\
Attributes & 1.00 & 80/20 ranks & $0.118/0.800$ & $0.092/0.783$ & $0.038/0.637$ \\
Graph & 0.25 & Zero/positive & $0.085/0.675$ & $0.089/0.778$ & $0.047/0.572$ \\
Graph & 0.25 & 80/20 ranks & $0.127/0.813$ & $0.089/0.778$ & $0.038/0.509$ \\
Graph & 0.50 & Zero/positive & $0.118/0.756$ & $0.091/0.791$ & $0.045/0.665$ \\
Graph & 0.50 & 80/20 ranks & $0.132/0.814$ & $0.091/0.791$ & $0.040/0.579$ \\
Graph & 1.00 & Zero/positive & $0.155/0.791$ & $0.091/0.793$ & $0.046/0.697$ \\
Graph & 1.00 & 80/20 ranks & $0.139/0.818$ & $0.091/0.793$ & $0.034/0.639$ \\
\bottomrule\end{tabular}\end{table*}

\begin{table*}[t]\centering\small
\caption{Tolokers: complete stratified comparison, mean FDP / power over ten training-seed averages. Both partitions and all revelation fractions are shown; no partition is selected using these outcomes.}
\begin{tabular}{lllrrr}\toprule Score & $\rho$ & Partition & Low reference & Full random & Stratified\\\midrule
Attributes & 0.25 & Zero/positive & $0.028/0.002$ & $0.009/0.000$ & $0.000/0.000$ \\
Attributes & 0.25 & 80/20 ranks & $0.000/0.000$ & $0.009/0.000$ & $0.000/0.000$ \\
Attributes & 0.50 & Zero/positive & $0.057/0.003$ & $0.000/0.000$ & $0.000/0.000$ \\
Attributes & 0.50 & 80/20 ranks & $0.000/0.000$ & $0.000/0.000$ & $0.000/0.000$ \\
Attributes & 1.00 & Zero/positive & $0.065/0.003$ & $0.000/0.000$ & $0.000/0.000$ \\
Attributes & 1.00 & 80/20 ranks & $0.000/0.000$ & $0.000/0.000$ & $0.000/0.000$ \\
Graph & 0.25 & Zero/positive & $0.000/0.000$ & $0.010/0.001$ & $0.000/0.000$ \\
Graph & 0.25 & 80/20 ranks & $0.006/0.001$ & $0.010/0.001$ & $0.000/0.000$ \\
Graph & 0.50 & Zero/positive & $0.005/0.001$ & $0.023/0.001$ & $0.000/0.000$ \\
Graph & 0.50 & 80/20 ranks & $0.040/0.002$ & $0.023/0.001$ & $0.000/0.000$ \\
Graph & 1.00 & Zero/positive & $0.005/0.001$ & $0.026/0.001$ & $0.004/0.000$ \\
Graph & 1.00 & 80/20 ranks & $0.044/0.002$ & $0.026/0.001$ & $0.000/0.000$ \\
\bottomrule\end{tabular}\end{table*}

\begin{table*}[t]\centering\small
\caption{Weibo (PyGOD): complete stratified comparison, mean FDP / power over ten training-seed averages. Both partitions and all revelation fractions are shown; no partition is selected using these outcomes.}
\begin{tabular}{lllrrr}\toprule Score & $\rho$ & Partition & Low reference & Full random & Stratified\\\midrule
Attributes & 0.25 & Zero/positive & $0.590/0.260$ & $0.000/0.000$ & $0.000/0.000$ \\
Attributes & 0.25 & 80/20 ranks & $0.592/0.262$ & $0.000/0.000$ & $0.000/0.000$ \\
Attributes & 0.50 & Zero/positive & $0.656/0.343$ & $0.013/0.001$ & $0.000/0.000$ \\
Attributes & 0.50 & 80/20 ranks & $0.656/0.342$ & $0.013/0.001$ & $0.000/0.000$ \\
Attributes & 1.00 & Zero/positive & $0.675/0.371$ & $0.015/0.000$ & $0.000/0.000$ \\
Attributes & 1.00 & 80/20 ranks & $0.676/0.373$ & $0.015/0.000$ & $0.000/0.000$ \\
Graph & 0.25 & Zero/positive & $0.666/0.629$ & $0.000/0.000$ & $0.000/0.000$ \\
Graph & 0.25 & 80/20 ranks & $0.666/0.630$ & $0.000/0.000$ & $0.000/0.000$ \\
Graph & 0.50 & Zero/positive & $0.683/0.768$ & $0.000/0.000$ & $0.000/0.000$ \\
Graph & 0.50 & 80/20 ranks & $0.684/0.774$ & $0.000/0.000$ & $0.000/0.000$ \\
Graph & 1.00 & Zero/positive & $0.695/0.826$ & $0.000/0.000$ & $0.000/0.000$ \\
Graph & 1.00 & 80/20 ranks & $0.697/0.843$ & $0.000/0.000$ & $0.000/0.000$ \\
\bottomrule\end{tabular}\end{table*}

\begin{table*}[t]\centering\small
\caption{T-Finance: complete stratified comparison, mean FDP / power over ten training-seed averages. Both partitions and all revelation fractions are shown; no partition is selected using these outcomes.}
\begin{tabular}{lllrrr}\toprule Score & $\rho$ & Partition & Low reference & Full random & Stratified\\\midrule
Attributes & 0.25 & Zero/positive & $0.067/0.627$ & $0.093/0.650$ & $0.046/0.633$ \\
Attributes & 0.25 & 80/20 ranks & $0.119/0.665$ & $0.093/0.650$ & $0.037/0.623$ \\
Attributes & 0.50 & Zero/positive & $0.067/0.626$ & $0.093/0.651$ & $0.045/0.633$ \\
Attributes & 0.50 & 80/20 ranks & $0.127/0.669$ & $0.093/0.651$ & $0.037/0.628$ \\
Attributes & 1.00 & Zero/positive & $0.071/0.633$ & $0.095/0.652$ & $0.045/0.629$ \\
Attributes & 1.00 & 80/20 ranks & $0.138/0.675$ & $0.095/0.652$ & $0.037/0.629$ \\
Graph & 0.25 & Zero/positive & $0.123/0.803$ & $0.092/0.787$ & $0.043/0.744$ \\
Graph & 0.25 & 80/20 ranks & $0.182/0.827$ & $0.092/0.787$ & $0.038/0.731$ \\
Graph & 0.50 & Zero/positive & $0.122/0.803$ & $0.095/0.789$ & $0.044/0.751$ \\
Graph & 0.50 & 80/20 ranks & $0.192/0.830$ & $0.095/0.789$ & $0.038/0.737$ \\
Graph & 1.00 & Zero/positive & $0.121/0.804$ & $0.099/0.792$ & $0.046/0.753$ \\
Graph & 1.00 & 80/20 ranks & $0.202/0.834$ & $0.099/0.792$ & $0.038/0.738$ \\
\bottomrule\end{tabular}\end{table*}

\begin{table*}[t]\centering\small
\caption{Weibo (GADBench): complete stratified comparison, mean FDP / power over ten training-seed averages. Both partitions and all revelation fractions are shown; no partition is selected using these outcomes.}
\begin{tabular}{lllrrr}\toprule Score & $\rho$ & Partition & Low reference & Full random & Stratified\\\midrule
Attributes & 0.25 & Zero/positive & $0.164/0.758$ & $0.093/0.664$ & $0.025/0.478$ \\
Attributes & 0.25 & 80/20 ranks & $0.168/0.764$ & $0.093/0.664$ & $0.017/0.401$ \\
Attributes & 0.50 & Zero/positive & $0.181/0.779$ & $0.096/0.666$ & $0.036/0.537$ \\
Attributes & 0.50 & 80/20 ranks & $0.189/0.787$ & $0.096/0.666$ & $0.025/0.489$ \\
Attributes & 1.00 & Zero/positive & $0.198/0.795$ & $0.093/0.664$ & $0.042/0.559$ \\
Attributes & 1.00 & 80/20 ranks & $0.216/0.811$ & $0.093/0.664$ & $0.026/0.495$ \\
Graph & 0.25 & Zero/positive & $0.253/0.966$ & $0.094/0.845$ & $0.030/0.571$ \\
Graph & 0.25 & 80/20 ranks & $0.258/0.969$ & $0.094/0.845$ & $0.023/0.474$ \\
Graph & 0.50 & Zero/positive & $0.280/0.972$ & $0.090/0.846$ & $0.029/0.630$ \\
Graph & 0.50 & 80/20 ranks & $0.290/0.974$ & $0.090/0.846$ & $0.020/0.523$ \\
Graph & 1.00 & Zero/positive & $0.311/0.978$ & $0.093/0.848$ & $0.035/0.673$ \\
Graph & 1.00 & 80/20 ranks & $0.317/0.978$ & $0.093/0.848$ & $0.023/0.567$ \\
\bottomrule\end{tabular}\end{table*}
